\documentclass[12pt,a4paper]{article}
\usepackage[utf8]{inputenc}\usepackage[T1]{fontenc}
\usepackage{geometry}\geometry{margin=2cm}
\usepackage{graphicx}\usepackage{amsmath,amssymb}\usepackage{float}
\usepackage{caption,subcaption}\usepackage{booktabs}\usepackage{hyperref}
\usepackage{listings}\usepackage{xcolor}\usepackage{parskip}
\lstdefinestyle{matlab}{language=Matlab,basicstyle=\ttfamily\scriptsize,
keywordstyle=\color{blue},commentstyle=\color{green!50!black},
stringstyle=\color{orange},numbers=left,numberstyle=\tiny\color{gray},
breaklines=true,frame=single,backgroundcolor=\color{gray!5}}
\title{\textbf{EEE 475/575 -- Medical Image Reconstruction}\\[0.3em]\large Homework \#2}
\author{\"Ozg\"ur Ekin S\"onmez}\date{31 March 2026}
\begin{document}\maketitle\tableofcontents\newpage

\section*{GenAI Usage Disclosure}
\addcontentsline{toc}{section}{GenAI Usage Disclosure}
I used Google Gemini (Antigravity) to help me structure the MATLAB code and format this report. I checked all results myself and I can explain every part.

\section{Part I: Gridding Reconstruction (Spiral Data)}

\subsection{K-Space Trajectory (1.1)}
The spiral k-space data consists of 6 interleaves with 2048 samples each. The trajectory spirals outward from the k-space center, giving denser sampling near DC and sparser sampling at higher spatial frequencies.
\begin{figure}[H]\centering\includegraphics[width=0.55\textwidth]{results/plots/1_1_spiral_trajectory.png}\caption{Spiral k-space trajectory.}\end{figure}
\begin{lstlisting}[style=matlab]
load('spiral.mat');
figure; plot(ktraj); xlabel('k_x'); ylabel('k_y');
title('Interleaved Spiral K-Space Trajectory'); axis equal; grid on;
saveas(gcf, fullfile(plotDir,'1_1_spiral_trajectory.png'));
\end{lstlisting}

\subsection{Density Compensation via Voronoi (1.2)}
The Voronoi tessellation assigns an area to each sample point. 186 samples at the boundary have NaN (unbounded cells). Some edge samples also have unreasonably large areas because their Voronoi cells extend far beyond the trajectory.
\begin{figure}[H]\centering
\begin{subfigure}{0.48\textwidth}\includegraphics[width=\textwidth]{results/plots/1_2_raw_area.png}\caption{Raw.}\end{subfigure}\hfill
\begin{subfigure}{0.48\textwidth}\includegraphics[width=\textwidth]{results/plots/1_2_raw_area_zoomed.png}\caption{Zoomed.}\end{subfigure}
\caption{Voronoi areas.}\end{figure}
\begin{lstlisting}[style=matlab]
area_raw = voronoidens(ktraj);
nan_count = sum(isnan(area_raw(:)));
fprintf('NaN count: %d\n', nan_count);
figure; plot(area_raw(:)); xlabel('Sample index'); ylabel('Area');
title('Raw Voronoi Area'); saveas(gcf, fullfile(plotDir,'1_2_raw_area.png'));
valid = area_raw(~isnan(area_raw(:)) & area_raw(:)<1e-2);
figure; plot(area_raw(:)); ylim([0 max(valid)*1.5]);
xlabel('Sample index'); ylabel('Area'); title('Voronoi Area - Zoomed');
saveas(gcf, fullfile(plotDir,'1_2_raw_area_zoomed.png'));
\end{lstlisting}

\subsection{Correcting the DCF (1.3)}
I replaced NaN and outlier values with a threshold of $3\times\text{median}$ of finite values ($\approx 0.000222$). This prevents boundary artifacts while preserving the correct weighting for interior samples.
\begin{figure}[H]\centering\includegraphics[width=0.5\textwidth]{results/plots/1_3_corrected_area.png}\caption{Corrected DCF.}\end{figure}
\begin{lstlisting}[style=matlab]
area_corr = area_raw;
fin = area_raw(isfinite(area_raw(:)));
thr = 3*median(fin);
area_corr(isnan(area_corr)) = thr;
area_corr(area_corr > thr)  = thr;
fprintf('Threshold: %.6f, remaining NaN: %d\n', thr, sum(isnan(area_corr(:))));
figure; plot(area_corr(:)); xlabel('Sample index'); ylabel('Area');
title('Corrected Area'); saveas(gcf, fullfile(plotDir,'1_3_corrected_area.png'));
\end{lstlisting}

\subsection{Direct Summation Reference (1.4)}
Implements the 2D NUFFT via direct summation: $m[n_x,n_y]=\sum_\ell d_\ell M(k_\ell) e^{j2\pi(k_{x,\ell}\Delta x(n_x-\tau)+k_{y,\ell}\Delta x(n_y-\tau))}$ with $N=128$, $\Delta x=1$, $\tau=64$. I loop row-by-row to keep memory low. This is the reference image for Part I. \textbf{CPU: 0.98 s.}
\begin{figure}[H]\centering\includegraphics[width=0.55\textwidth]{results/plots/1_4_image.png}\caption{Reference image.}\end{figure}
\begin{figure}[H]\centering
\begin{subfigure}{0.48\textwidth}\includegraphics[width=\textwidth]{results/plots/1_4_horiz.png}\caption{Horizontal.}\end{subfigure}\hfill
\begin{subfigure}{0.48\textwidth}\includegraphics[width=\textwidth]{results/plots/1_4_vert.png}\caption{Vertical.}\end{subfigure}
\caption{Cross-sections.}\end{figure}
\begin{lstlisting}[style=matlab]
N=128; dx=1; tau=N/2;
kx_vec=real(ktraj(:)); ky_vec=imag(ktraj(:));
M_vec=kdata(:); d_vec=area_corr(:);
dM = d_vec.*M_vec;
n_vals = (0:N-1)-tau;
phase_y = exp(1j*2*pi*dx * ky_vec * n_vals);
ima_direct = zeros(N,N);
t0 = cputime;
for ix = 1:N
    phase_x = exp(1j*2*pi*dx * kx_vec * n_vals(ix));
    ima_direct(ix,:) = (dM.*phase_x).' * phase_y;
end
t_direct = cputime-t0; clear phase_y phase_x;
ima_direct = normImg(ima_direct.');
fprintf('1.4 CPU time: %.2f s\n', t_direct);
figure; imshow(ima_direct,[0 1]); colormap gray; colorbar;
title('Direct Summation 128x128'); saveas(gcf, fullfile(plotDir,'1_4_image.png'));
figure; plot(abs(ima_direct(end/2,:))); xlabel('Column'); ylabel('Intensity');
title('1.4 Horizontal Cross-Section'); grid on; saveas(gcf,fullfile(plotDir,'1_4_horiz.png'));
figure; plot(abs(ima_direct(:,end/2))); xlabel('Row'); ylabel('Intensity');
title('1.4 Vertical Cross-Section'); grid on; saveas(gcf,fullfile(plotDir,'1_4_vert.png'));
save(fullfile(matDir,'ima_direct.mat'),'ima_direct');
\end{lstlisting}

\subsection{No Density Compensation (1.5)}
Without DCF, low frequencies are overweighted because the spiral samples more densely near the k-space center. The image appears blurred and washed out. \textbf{PSNR=10.86 dB, SSIM=0.2729, CPU=0.95 s.} This confirms that density compensation is essential for non-Cartesian reconstruction.
\begin{figure}[H]\centering
\begin{subfigure}{0.48\textwidth}\includegraphics[width=\textwidth]{results/plots/1_5_image.png}\caption{Image.}\end{subfigure}\hfill
\begin{subfigure}{0.48\textwidth}\includegraphics[width=\textwidth]{results/plots/1_5_error.png}\caption{Error.}\end{subfigure}
\caption{1.5 No DCF.}\end{figure}
\begin{figure}[H]\centering
\begin{subfigure}{0.48\textwidth}\includegraphics[width=\textwidth]{results/plots/1_5_horiz.png}\caption{Horiz XS.}\end{subfigure}\hfill
\begin{subfigure}{0.48\textwidth}\includegraphics[width=\textwidth]{results/plots/1_5_vert.png}\caption{Vert XS.}\end{subfigure}
\caption{1.5 Cross-sections.}\end{figure}
\begin{lstlisting}[style=matlab]
n_vals = (0:N-1)-tau;
phase_y = exp(1j*2*pi*dx * ky_vec * n_vals);
ima_nodcf = zeros(N,N);
t0 = cputime;
for ix = 1:N
    phase_x = exp(1j*2*pi*dx * kx_vec * n_vals(ix));
    ima_nodcf(ix,:) = (M_vec.*phase_x).' * phase_y;
end
t_nodcf = cputime-t0; clear phase_y phase_x;
ima_nodcf = normImg(ima_nodcf.');
err_nodcf = abs(ima_nodcf - ima_direct);
psnr_nodcf = psnr(ima_nodcf, ima_direct);
ssim_nodcf = ssim(ima_nodcf, ima_direct);
fprintf('1.5 No DCF: PSNR=%.2f dB, SSIM=%.4f, CPU=%.2f s\n',psnr_nodcf,ssim_nodcf,t_nodcf);
figure; imshow(ima_nodcf,[0 1]); colormap gray; colorbar;
title('1.5 No DCF Image'); saveas(gcf, fullfile(plotDir,'1_5_image.png'));
figure; imshow(err_nodcf,[0 0.2]); colormap gray; colorbar;
title('1.5 Error [0 0.2]'); saveas(gcf, fullfile(plotDir,'1_5_error.png'));
figure; plot(abs(ima_nodcf(end/2,:)),'b'); hold on;
plot(abs(ima_direct(end/2,:)),'r--'); legend('NoDCF','Ref');
xlabel('Column'); title('1.5 Horizontal XS'); grid on; saveas(gcf,fullfile(plotDir,'1_5_horiz.png'));
figure; plot(abs(ima_nodcf(:,end/2)),'b'); hold on;
plot(abs(ima_direct(:,end/2)),'r--'); legend('NoDCF','Ref');
xlabel('Row'); title('1.5 Vertical XS'); grid on; saveas(gcf,fullfile(plotDir,'1_5_vert.png'));
\end{lstlisting}

\subsection{Double FOV (1.6)}
With $N=256$, $\Delta x=1$, FOV doubles. The object appears in the center with empty space around it. The cropped $128\times128$ center is identical to the reference (PSNR=$\infty$, SSIM=1.0000). No new information is gained but computation took 2.36 s vs 0.98 s. \textbf{CPU=2.36 s.}
\begin{figure}[H]\centering
\begin{subfigure}{0.48\textwidth}\includegraphics[width=\textwidth]{results/plots/1_6_full.png}\caption{Full 256x256.}\end{subfigure}\hfill
\begin{subfigure}{0.48\textwidth}\includegraphics[width=\textwidth]{results/plots/1_6_crop.png}\caption{Cropped 128x128.}\end{subfigure}
\caption{1.6 Double FOV.}\end{figure}
\begin{figure}[H]\centering
\begin{subfigure}{0.48\textwidth}\includegraphics[width=\textwidth]{results/plots/1_6_error.png}\caption{Error.}\end{subfigure}\hfill
\begin{subfigure}{0.48\textwidth}\includegraphics[width=\textwidth]{results/plots/1_6_horiz.png}\caption{Horizontal XS.}\end{subfigure}
\caption{1.6 Error and horizontal cross-section.}\end{figure}
\begin{figure}[H]\centering
\includegraphics[width=0.48\textwidth]{results/plots/1_6_vert.png}
\caption{1.6 Vertical cross-section.}\end{figure}
\begin{lstlisting}[style=matlab]
N2=256; tau2=N2/2; n2=(0:N2-1)-tau2;
phase_y2 = exp(1j*2*pi*dx * ky_vec * n2);
ima_dfov = zeros(N2,N2);
t0 = cputime;
for ix = 1:N2
    px = exp(1j*2*pi*dx * kx_vec * n2(ix));
    ima_dfov(ix,:) = (dM.*px).' * phase_y2;
end
t_dfov = cputime-t0; clear phase_y2 px;
ima_dfov = normImg(ima_dfov.');
c=N2/2; h=N/2;
ima_dfov_crop = normImg(ima_dfov(c-h+1:c+h, c-h+1:c+h));
err_dfov = abs(ima_dfov_crop - ima_direct);
psnr_dfov=psnr(ima_dfov_crop,ima_direct); ssim_dfov=ssim(ima_dfov_crop,ima_direct);
fprintf('1.6: PSNR=%.2f, SSIM=%.4f, CPU=%.2f\n',psnr_dfov,ssim_dfov,t_dfov);
figure; imshow(ima_dfov,[0 1]); colormap gray; colorbar;
title('1.6 Full 256x256'); saveas(gcf, fullfile(plotDir,'1_6_full.png'));
figure; imshow(ima_dfov_crop,[0 1]); colormap gray; colorbar;
title('1.6 Cropped 128x128'); saveas(gcf, fullfile(plotDir,'1_6_crop.png'));
figure; imshow(err_dfov,[0 0.2]); colormap gray; colorbar;
title('1.6 Error'); saveas(gcf, fullfile(plotDir,'1_6_error.png'));
figure; plot(abs(ima_dfov_crop(end/2,:)),'b'); hold on;
plot(abs(ima_direct(end/2,:)),'r--'); legend('DblFOV','Ref');
title('1.6 Horiz XS'); grid on; saveas(gcf, fullfile(plotDir,'1_6_horiz.png'));
figure; plot(abs(ima_dfov_crop(:,end/2)),'b'); hold on;
plot(abs(ima_direct(:,end/2)),'r--'); legend('DblFOV','Ref');
title('1.6 Vert XS'); grid on; saveas(gcf, fullfile(plotDir,'1_6_vert.png'));
\end{lstlisting}

\subsection{Half Pixel Size (1.7)}
With $\Delta x=0.5$, $N=256$, same FOV but finer pixels. Reference is zero-padded in k-space to $256\times256$ for fair comparison. The result is very close to the reference. \textbf{PSNR=33.67 dB, SSIM=0.9150, CPU=3.62 s.} Higher resolution image without new frequency content.
\begin{figure}[H]\centering
\begin{subfigure}{0.32\textwidth}\includegraphics[width=\textwidth]{results/plots/1_7_image.png}\caption{Image.}\end{subfigure}\hfill
\begin{subfigure}{0.32\textwidth}\includegraphics[width=\textwidth]{results/plots/1_7_ref.png}\caption{Zero-padded ref.}\end{subfigure}\hfill
\begin{subfigure}{0.32\textwidth}\includegraphics[width=\textwidth]{results/plots/1_7_error.png}\caption{Error.}\end{subfigure}
\caption{1.7 Half pixel.}\end{figure}
\begin{figure}[H]\centering
\begin{subfigure}{0.48\textwidth}\includegraphics[width=\textwidth]{results/plots/1_7_horiz.png}\caption{Horizontal XS.}\end{subfigure}\hfill
\begin{subfigure}{0.48\textwidth}\includegraphics[width=\textwidth]{results/plots/1_7_vert.png}\caption{Vertical XS.}\end{subfigure}
\caption{1.7 Cross-sections.}\end{figure}
\begin{lstlisting}[style=matlab]
N3=256; dx3=0.5; tau3=N3/2; n3=(0:N3-1)-tau3;
phase_y3 = exp(1j*2*pi*dx3 * ky_vec * n3);
ima_hpx = zeros(N3,N3);
t0=cputime;
for ix=1:N3
    px=exp(1j*2*pi*dx3*kx_vec*n3(ix));
    ima_hpx(ix,:)=(dM.*px).'*phase_y3;
end
t_hpx=cputime-t0; clear phase_y3 px;
ima_hpx = normImg(ima_hpx.');
kref=fftshift(fft2(ifftshift(ima_direct)));
pad=(N3-N)/2; kref_pad=zeros(N3,N3);
kref_pad(pad+1:pad+N, pad+1:pad+N)=kref;
ima_ref256=normImg(fftshift(ifft2(ifftshift(kref_pad))));
err_hpx=abs(ima_hpx-ima_ref256);
psnr_hpx=psnr(ima_hpx,ima_ref256); ssim_hpx=ssim(ima_hpx,ima_ref256);
figure; imshow(ima_hpx,[0 1]); colormap gray; colorbar;
title('1.7 Half Pixel'); saveas(gcf, fullfile(plotDir,'1_7_image.png'));
figure; imshow(err_hpx,[0 0.2]); colormap gray; colorbar;
title('1.7 Error'); saveas(gcf, fullfile(plotDir,'1_7_error.png'));
\end{lstlisting}

\subsection{Double Pixel Size (1.8)}
With $\Delta x=2$, $N=64$, same FOV but coarser pixels. Result is zero-padded from $64\times64$ to $128\times128$ for comparison. Fine details are visibly lost. \textbf{PSNR=18.95 dB, SSIM=0.6160, CPU=0.36 s.} Faster but too much quality loss.
\begin{figure}[H]\centering
\begin{subfigure}{0.32\textwidth}\includegraphics[width=\textwidth]{results/plots/1_8_image.png}\caption{64x64.}\end{subfigure}\hfill
\begin{subfigure}{0.32\textwidth}\includegraphics[width=\textwidth]{results/plots/1_8_zp.png}\caption{Zero-padded.}\end{subfigure}\hfill
\begin{subfigure}{0.32\textwidth}\includegraphics[width=\textwidth]{results/plots/1_8_error.png}\caption{Error.}\end{subfigure}
\caption{1.8 Double pixel.}\end{figure}
\begin{figure}[H]\centering
\begin{subfigure}{0.48\textwidth}\includegraphics[width=\textwidth]{results/plots/1_8_horiz.png}\caption{Horizontal XS.}\end{subfigure}\hfill
\begin{subfigure}{0.48\textwidth}\includegraphics[width=\textwidth]{results/plots/1_8_vert.png}\caption{Vertical XS.}\end{subfigure}
\caption{1.8 Cross-sections.}\end{figure}
\begin{lstlisting}[style=matlab]
N4=64; dx4=2; tau4=N4/2; n4=(0:N4-1)-tau4;
phase_y4=exp(1j*2*pi*dx4*ky_vec*n4);
ima_dpx=zeros(N4,N4);
t0=cputime;
for ix=1:N4
    px=exp(1j*2*pi*dx4*kx_vec*n4(ix));
    ima_dpx(ix,:)=(dM.*px).'*phase_y4;
end
t_dpx=cputime-t0; clear phase_y4 px;
ima_dpx=normImg(ima_dpx.');
kdpx=fftshift(fft2(ifftshift(ima_dpx)));
pad2=(N-N4)/2; kdpx_pad=zeros(N,N);
kdpx_pad(pad2+1:pad2+N4,pad2+1:pad2+N4)=kdpx;
ima_dpx_zp=normImg(fftshift(ifft2(ifftshift(kdpx_pad))));
err_dpx=abs(ima_dpx_zp-ima_direct);
psnr_dpx=psnr(ima_dpx_zp,ima_direct); ssim_dpx=ssim(ima_dpx_zp,ima_direct);
figure; imshow(ima_dpx,[0 1]); colormap gray; colorbar;
title('1.8 Original 64x64'); saveas(gcf, fullfile(plotDir,'1_8_image.png'));
figure; imshow(err_dpx,[0 0.2]); colormap gray; colorbar;
title('1.8 Error'); saveas(gcf, fullfile(plotDir,'1_8_error.png'));
\end{lstlisting}

\subsection{1X Gridding (1.9)}
Using \texttt{gridkb} with osf=1, wg=2. I transpose the gridding output to match the image-axis convention used by the direct-summation reference. The result is much closer to the reference than before, but some residual artifacts remain because the grid is not oversampled. \textbf{PSNR=23.91 dB, SSIM=0.7537, CPU=0.05 s.}
\begin{figure}[H]\centering
\begin{subfigure}{0.48\textwidth}\includegraphics[width=\textwidth]{results/plots/1_9_image.png}\caption{Image.}\end{subfigure}\hfill
\begin{subfigure}{0.48\textwidth}\includegraphics[width=\textwidth]{results/plots/1_9_error.png}\caption{Error.}\end{subfigure}
\caption{1.9 Gridding 1X.}\end{figure}
\begin{figure}[H]\centering
\begin{subfigure}{0.48\textwidth}\includegraphics[width=\textwidth]{results/plots/1_9_horiz.png}\caption{Horizontal XS.}\end{subfigure}\hfill
\begin{subfigure}{0.48\textwidth}\includegraphics[width=\textwidth]{results/plots/1_9_vert.png}\caption{Vertical XS.}\end{subfigure}
\caption{1.9 Cross-sections.}\end{figure}
\begin{lstlisting}[style=matlab]
t0=cputime;
ima_g1=gridkb(kdata,ktraj,area_corr,128,1,2,'image');
t_g1=cputime-t0;
ima_g1=normImg(ima_g1.');
err_g1=abs(ima_g1-ima_direct);
psnr_g1=psnr(ima_g1,ima_direct); ssim_g1=ssim(ima_g1,ima_direct);
figure; imshow(ima_g1,[0 1]); colormap gray; colorbar;
title('1.9 Gridding 1X'); saveas(gcf, fullfile(plotDir,'1_9_image.png'));
figure; imshow(err_g1,[0 0.2]); colormap gray; colorbar;
title('1.9 Error'); saveas(gcf, fullfile(plotDir,'1_9_error.png'));
figure; plot(abs(ima_g1(end/2,:)),'b'); hold on;
plot(abs(ima_direct(end/2,:)),'r--'); legend('1XGrid','Ref');
title('1.9 Horiz XS'); grid on; saveas(gcf, fullfile(plotDir,'1_9_horiz.png'));
figure; plot(abs(ima_g1(:,end/2)),'b'); hold on;
plot(abs(ima_direct(:,end/2)),'r--'); legend('1XGrid','Ref');
title('1.9 Vert XS'); grid on; saveas(gcf, fullfile(plotDir,'1_9_vert.png'));
\end{lstlisting}

\subsection{2X Gridding (1.10)}
With osf=2, wg=4. Grid is $256\times256$; I transpose the gridding output to match the reference orientation, then crop the central $128\times128$. The wider kernel and oversampling reduce aliasing strongly, and the result becomes nearly identical to the reference. \textbf{PSNR=44.35 dB, SSIM=0.9906, CPU=0.09 s.} This is a major improvement over 1X gridding.
\begin{figure}[H]\centering
\begin{subfigure}{0.32\textwidth}\includegraphics[width=\textwidth]{results/plots/1_10_full.png}\caption{Full 256x256.}\end{subfigure}\hfill
\begin{subfigure}{0.32\textwidth}\includegraphics[width=\textwidth]{results/plots/1_10_crop.png}\caption{Cropped.}\end{subfigure}\hfill
\begin{subfigure}{0.32\textwidth}\includegraphics[width=\textwidth]{results/plots/1_10_error.png}\caption{Error.}\end{subfigure}
\caption{1.10 Gridding 2X.}\end{figure}
\begin{figure}[H]\centering
\begin{subfigure}{0.48\textwidth}\includegraphics[width=\textwidth]{results/plots/1_10_horiz.png}\caption{Horizontal XS.}\end{subfigure}\hfill
\begin{subfigure}{0.48\textwidth}\includegraphics[width=\textwidth]{results/plots/1_10_vert.png}\caption{Vertical XS.}\end{subfigure}
\caption{1.10 Cross-sections.}\end{figure}
\begin{lstlisting}[style=matlab]
t0=cputime;
ima_g2_full=gridkb(kdata,ktraj,area_corr,128,2,4,'image');
t_g2=cputime-t0;
ima_g2_full=normImg(ima_g2_full.');
c2=128; h2=64;
ima_g2=normImg(ima_g2_full(c2-h2+1:c2+h2,c2-h2+1:c2+h2));
err_g2=abs(ima_g2-ima_direct);
psnr_g2=psnr(ima_g2,ima_direct); ssim_g2=ssim(ima_g2,ima_direct);
figure; imshow(ima_g2_full,[0 1]); colormap gray; colorbar;
title('1.10 Full 256x256'); saveas(gcf, fullfile(plotDir,'1_10_full.png'));
figure; imshow(ima_g2,[0 1]); colormap gray; colorbar;
title('1.10 Cropped 128x128'); saveas(gcf, fullfile(plotDir,'1_10_crop.png'));
figure; imshow(err_g2,[0 0.2]); colormap gray; colorbar;
title('1.10 Error'); saveas(gcf, fullfile(plotDir,'1_10_error.png'));
figure; plot(abs(ima_g2(end/2,:)),'b'); hold on;
plot(abs(ima_direct(end/2,:)),'r--'); legend('2XGrid','Ref');
title('1.10 Horiz XS'); grid on; saveas(gcf, fullfile(plotDir,'1_10_horiz.png'));
figure; plot(abs(ima_g2(:,end/2)),'b'); hold on;
plot(abs(ima_direct(:,end/2)),'r--'); legend('2XGrid','Ref');
title('1.10 Vert XS'); grid on; saveas(gcf, fullfile(plotDir,'1_10_vert.png'));
\end{lstlisting}

\subsection{Scattered Interpolation 1X (1.11)}
Using \texttt{scatteredInterpolant} with linear interpolation on a $128\times128$ grid from $\pm0.5$. NaN values (outside the spiral ROS) are set to zero. This uses only nearby data points (unlike gridding which convolves with all). The image may be flipped/transposed; I corrected the orientation before comparison. \textbf{PSNR=22.85 dB, SSIM=0.6705, CPU=0.06 s.} It gives a reasonable reconstruction, but after fixing the gridding orientation it is now worse than both 1X and especially 2X gridding for this dataset.
\begin{figure}[H]\centering
\begin{subfigure}{0.48\textwidth}\includegraphics[width=\textwidth]{results/plots/1_11_image.png}\caption{Image.}\end{subfigure}\hfill
\begin{subfigure}{0.48\textwidth}\includegraphics[width=\textwidth]{results/plots/1_11_error.png}\caption{Error.}\end{subfigure}
\caption{1.11 ScatInterp 1X.}\end{figure}
\begin{figure}[H]\centering
\begin{subfigure}{0.48\textwidth}\includegraphics[width=\textwidth]{results/plots/1_11_horiz.png}\caption{Horizontal XS.}\end{subfigure}\hfill
\begin{subfigure}{0.48\textwidth}\includegraphics[width=\textwidth]{results/plots/1_11_vert.png}\caption{Vertical XS.}\end{subfigure}
\caption{1.11 Cross-sections.}\end{figure}
\begin{lstlisting}[style=matlab]
kx_s=real(ktraj(:)); ky_s=imag(ktraj(:));
F_re = scatteredInterpolant(kx_s,ky_s,real(kdata(:)),'linear','none');
F_im = scatteredInterpolant(kx_s,ky_s,imag(kdata(:)),'linear','none');
kg1d = linspace(-0.5,0.5,128);
[KX1,KY1] = meshgrid(kg1d,kg1d);
t0=cputime;
kgrid1 = F_re(KX1,KY1) + 1j*F_im(KX1,KY1);
t_s1=cputime-t0;
kgrid1(isnan(kgrid1))=0;
ima_s1 = fftshift(ifft2(ifftshift(kgrid1)));
ima_s1 = normImg(ima_s1);
candidates = {ima_s1, fliplr(ima_s1), flipud(ima_s1), ima_s1.', ...
              flipud(fliplr(ima_s1)), rot90(ima_s1,1), rot90(ima_s1,3)};
best_p=-Inf;
for ci=1:numel(candidates)
    if all(size(candidates{ci})==[N N])
        pp=psnr(candidates{ci},ima_direct);
        if pp>best_p, best_p=pp; ima_s1=candidates{ci}; end
    end
end
err_s1=abs(ima_s1-ima_direct);
psnr_s1=psnr(ima_s1,ima_direct); ssim_s1=ssim(ima_s1,ima_direct);
figure; imshow(ima_s1,[0 1]); colormap gray; colorbar;
title('1.11 ScatInterp 1X'); saveas(gcf, fullfile(plotDir,'1_11_image.png'));
figure; imshow(err_s1,[0 0.2]); colormap gray; colorbar;
title('1.11 Error'); saveas(gcf, fullfile(plotDir,'1_11_error.png'));
\end{lstlisting}

\subsection{Scattered Interpolation 2X (1.12)}
Same as 1.11 but with a $256\times256$ grid. Cropped central $128\times128$. The finer grid improves interpolation. \textbf{PSNR=24.33 dB, SSIM=0.7350, CPU=0.27 s.} It improves over 1X scattered interpolation, but it still remains below the oriented 2X gridding result for this dataset.
\begin{figure}[H]\centering
\begin{subfigure}{0.32\textwidth}\includegraphics[width=\textwidth]{results/plots/1_12_full.png}\caption{Full 256x256.}\end{subfigure}\hfill
\begin{subfigure}{0.32\textwidth}\includegraphics[width=\textwidth]{results/plots/1_12_crop.png}\caption{Cropped.}\end{subfigure}\hfill
\begin{subfigure}{0.32\textwidth}\includegraphics[width=\textwidth]{results/plots/1_12_error.png}\caption{Error.}\end{subfigure}
\caption{1.12 ScatInterp 2X.}\end{figure}
\begin{figure}[H]\centering
\begin{subfigure}{0.48\textwidth}\includegraphics[width=\textwidth]{results/plots/1_12_horiz.png}\caption{Horizontal XS.}\end{subfigure}\hfill
\begin{subfigure}{0.48\textwidth}\includegraphics[width=\textwidth]{results/plots/1_12_vert.png}\caption{Vertical XS.}\end{subfigure}
\caption{1.12 Cross-sections.}\end{figure}
\begin{lstlisting}[style=matlab]
kg2d = linspace(-0.5,0.5,256);
[KX2,KY2] = meshgrid(kg2d,kg2d);
t0=cputime;
kgrid2 = F_re(KX2,KY2) + 1j*F_im(KX2,KY2);
t_s2=cputime-t0;
kgrid2(isnan(kgrid2))=0;
ima_s2_full = fftshift(ifft2(ifftshift(kgrid2)));
ima_s2_full = normImg(ima_s2_full);
candidates2 = {ima_s2_full, fliplr(ima_s2_full), flipud(ima_s2_full), ...
    ima_s2_full.', flipud(fliplr(ima_s2_full)), rot90(ima_s2_full,1), rot90(ima_s2_full,3)};
best_p=-Inf;
for ci=1:numel(candidates2)
    tmp = candidates2{ci};
    if all(size(tmp)==[256 256])
        c3=128; crop=normImg(tmp(c3-h2+1:c3+h2,c3-h2+1:c3+h2));
        pp=psnr(crop,ima_direct);
        if pp>best_p, best_p=pp; ima_s2_full=tmp; ima_s2=crop; end
    end
end
err_s2=abs(ima_s2-ima_direct);
psnr_s2=psnr(ima_s2,ima_direct); ssim_s2=ssim(ima_s2,ima_direct);
figure; imshow(ima_s2_full,[0 1]); colormap gray; colorbar;
title('1.12 Full 256x256'); saveas(gcf, fullfile(plotDir,'1_12_full.png'));
figure; imshow(ima_s2,[0 1]); colormap gray; colorbar;
title('1.12 Cropped'); saveas(gcf, fullfile(plotDir,'1_12_crop.png'));
figure; imshow(err_s2,[0 0.2]); colormap gray; colorbar;
title('1.12 Error'); saveas(gcf, fullfile(plotDir,'1_12_error.png'));
\end{lstlisting}


\section{Part II: Gridding vs Backprojection}

\subsection{Phantom and Sinogram (2.1)}
Generated a $256\times256$ Modified Shepp-Logan phantom and computed its Radon transform for $0^\circ$--$179^\circ$. The sinogram is $367\times180$ where the $l$-axis shows the projection samples and $\theta$ shows the angle.
\begin{figure}[H]\centering
\begin{subfigure}{0.45\textwidth}\includegraphics[width=\textwidth]{results/plots/2_1_phantom.png}\caption{Phantom.}\end{subfigure}\hfill
\begin{subfigure}{0.50\textwidth}\includegraphics[width=\textwidth]{results/plots/2_1_sinogram.png}\caption{Sinogram.}\end{subfigure}
\caption{Phantom and sinogram.}\end{figure}
\begin{figure}[H]\centering
\begin{subfigure}{0.48\textwidth}\includegraphics[width=\textwidth]{results/plots/2_1_phantom_horiz.png}\caption{Horizontal XS.}\end{subfigure}\hfill
\begin{subfigure}{0.48\textwidth}\includegraphics[width=\textwidth]{results/plots/2_1_phantom_vert.png}\caption{Vertical XS.}\end{subfigure}
\caption{2.1 Phantom cross-sections.}\end{figure}
\begin{lstlisting}[style=matlab]
P = phantom('Modified Shepp-Logan',256);
P = normImg(P);
angles = 0:179;
proj = radon(P, angles);
figure; imshow(P,[0 1]); colormap gray; colorbar;
title('Shepp-Logan Phantom 256x256'); saveas(gcf, fullfile(plotDir,'2_1_phantom.png'));
figure; imagesc(angles,1:size(proj,1),proj); colormap gray; colorbar;
xlabel('\theta (degrees)'); ylabel('l'); title('Sinogram'); axis xy;
saveas(gcf, fullfile(plotDir,'2_1_sinogram.png'));
figure; plot(P(128,:)); xlabel('Column'); title('2.1 Phantom Horiz XS'); grid on;
saveas(gcf, fullfile(plotDir,'2_1_phantom_horiz.png'));
figure; plot(P(:,128)); xlabel('Row'); title('2.1 Phantom Vert XS'); grid on;
saveas(gcf, fullfile(plotDir,'2_1_phantom_vert.png'));
\end{lstlisting}

\subsection{Filtered Backprojection (2.2)}
Used \texttt{iradon} with the default Ram-Lak (ramp) filter. The ramp filter corrects the $1/|f|$ oversampling at low frequencies that occurs in radial sampling. Cropped the result to $256\times256$. \textbf{PSNR=26.82 dB, SSIM=0.6799, CPU=0.08 s.} Good reconstruction with minor streak artifacts from limited projections (180).
\begin{figure}[H]\centering
\begin{subfigure}{0.48\textwidth}\includegraphics[width=\textwidth]{results/plots/2_2_image.png}\caption{FBP.}\end{subfigure}\hfill
\begin{subfigure}{0.48\textwidth}\includegraphics[width=\textwidth]{results/plots/2_2_error.png}\caption{Error.}\end{subfigure}
\caption{2.2 FBP.}\end{figure}
\begin{figure}[H]\centering
\begin{subfigure}{0.48\textwidth}\includegraphics[width=\textwidth]{results/plots/2_2_horiz.png}\caption{Horizontal XS.}\end{subfigure}\hfill
\begin{subfigure}{0.48\textwidth}\includegraphics[width=\textwidth]{results/plots/2_2_vert.png}\caption{Vertical XS.}\end{subfigure}
\caption{2.2 Cross-sections.}\end{figure}
\begin{lstlisting}[style=matlab]
t0=cputime;
ima_fbp_raw = iradon(proj, angles);
t_fbp=cputime-t0;
csz=round(size(ima_fbp_raw,1)/2);
ima_fbp=ima_fbp_raw(csz-127:csz+128, csz-127:csz+128);
ima_fbp=normImg(ima_fbp);
err_fbp=abs(ima_fbp-P);
psnr_fbp=psnr(ima_fbp,P); ssim_fbp=ssim(ima_fbp,P);
figure; imshow(ima_fbp,[0 1]); colormap gray; colorbar;
title('2.2 FBP Image'); saveas(gcf, fullfile(plotDir,'2_2_image.png'));
figure; imshow(err_fbp,[0 0.2]); colormap gray; colorbar;
title('2.2 Error'); saveas(gcf, fullfile(plotDir,'2_2_error.png'));
figure; plot(ima_fbp(128,:),'b'); hold on; plot(P(128,:),'r--');
legend('FBP','Ref'); title('2.2 Horiz XS'); grid on; saveas(gcf,fullfile(plotDir,'2_2_horiz.png'));
figure; plot(ima_fbp(:,128),'b'); hold on; plot(P(:,128),'r--');
legend('FBP','Ref'); title('2.2 Vert XS'); grid on; saveas(gcf,fullfile(plotDir,'2_2_vert.png'));
\end{lstlisting}

\subsection{Naive Backprojection (2.3)}
Used \texttt{iradon} with filter set to \texttt{'none'}. Without the ramp filter, low frequencies dominate completely and the image is severely blurred. \textbf{PSNR=4.72 dB, SSIM=0.1535, CPU=0.08 s.} This shows the ramp filter is essential for backprojection.
\begin{figure}[H]\centering
\begin{subfigure}{0.48\textwidth}\includegraphics[width=\textwidth]{results/plots/2_3_image.png}\caption{Naive BP.}\end{subfigure}\hfill
\begin{subfigure}{0.48\textwidth}\includegraphics[width=\textwidth]{results/plots/2_3_error.png}\caption{Error.}\end{subfigure}
\caption{2.3 Naive backprojection.}\end{figure}
\begin{figure}[H]\centering
\begin{subfigure}{0.48\textwidth}\includegraphics[width=\textwidth]{results/plots/2_3_horiz.png}\caption{Horizontal XS.}\end{subfigure}\hfill
\begin{subfigure}{0.48\textwidth}\includegraphics[width=\textwidth]{results/plots/2_3_vert.png}\caption{Vertical XS.}\end{subfigure}
\caption{2.3 Cross-sections.}\end{figure}
\begin{lstlisting}[style=matlab]
t0=cputime;
ima_bp_raw = iradon(proj, angles, 'linear', 'none');
t_bp=cputime-t0;
ima_bp=ima_bp_raw(csz-127:csz+128, csz-127:csz+128);
ima_bp=normImg(ima_bp);
err_bp=abs(ima_bp-P);
psnr_bp=psnr(ima_bp,P); ssim_bp=ssim(ima_bp,P);
figure; imshow(ima_bp,[0 1]); colormap gray; colorbar;
title('2.3 Naive BP'); saveas(gcf, fullfile(plotDir,'2_3_image.png'));
figure; imshow(err_bp,[0 0.2]); colormap gray; colorbar;
title('2.3 Error'); saveas(gcf, fullfile(plotDir,'2_3_error.png'));
figure; plot(ima_bp(128,:),'b'); hold on; plot(P(128,:),'r--');
legend('NaiveBP','Ref'); title('2.3 Horiz XS'); grid on; saveas(gcf,fullfile(plotDir,'2_3_horiz.png'));
figure; plot(ima_bp(:,128),'b'); hold on; plot(P(:,128),'r--');
legend('NaiveBP','Ref'); title('2.3 Vert XS'); grid on; saveas(gcf,fullfile(plotDir,'2_3_vert.png'));
\end{lstlisting}

\subsection{K-Space via Projection-Slice Theorem (2.4)}
The Projection-Slice Theorem says the 1D FT of a projection at angle $\theta$ equals a radial line through 2D k-space at that angle. I set $k_{max}=0.5$ for compatibility with \texttt{gridkb}. The k-space trajectory is: $(k_x,k_y) = k_r(\cos\theta, \sin\theta)$. The log-magnitude plot shows the energy concentrated at the center.
\begin{figure}[H]\centering
\begin{subfigure}{0.48\textwidth}\includegraphics[width=\textwidth]{results/plots/2_4_trajectory.png}\caption{Trajectory.}\end{subfigure}\hfill
\begin{subfigure}{0.48\textwidth}\includegraphics[width=\textwidth]{results/plots/2_4_kspace_mag.png}\caption{K-space (log).}\end{subfigure}
\caption{2.4 Projection-Slice Theorem.}\end{figure}
\begin{lstlisting}[style=matlab]
Np = size(proj,1); kmax = 0.5;
kr_vec = linspace(-kmax, kmax, Np).';
kdata_proj = zeros(Np,180);
for th=1:180
    kdata_proj(:,th) = fftshift(fft(ifftshift(proj(:,th))));
end
theta_rad = deg2rad(angles);
kx_proj = kr_vec * cos(theta_rad);
ky_proj = kr_vec * sin(theta_rad);
ktraj_proj = kx_proj + 1j*ky_proj;
figure; plot(real(ktraj_proj(:)),imag(ktraj_proj(:)),'b.','MarkerSize',1);
xlabel('k_x'); ylabel('k_y'); title('2.4 Radial Trajectory'); axis equal; grid on;
saveas(gcf, fullfile(plotDir,'2_4_trajectory.png'));
figure; imagesc(angles, kr_vec, log(1+abs(kdata_proj))); colormap gray; colorbar;
xlabel('\theta'); ylabel('k_r'); title('2.4 K-Space Magnitude (log)'); axis xy;
saveas(gcf, fullfile(plotDir,'2_4_kspace_mag.png'));
\end{lstlisting}

\subsection{Geometric DCF (2.5)}
For radial trajectories, the area associated with each sample is $w(k_r)=|k_r|\cdot\frac{2\pi}{N_\theta}\cdot\Delta k_r$. The DC point gets a small disk area $\pi(\Delta k_r/2)^2/N_\theta$. I cannot use \texttt{voronoidens} here because multiple lines pass through the center, giving duplicate positions.
\begin{figure}[H]\centering\includegraphics[width=0.5\textwidth]{results/plots/2_5_dcf.png}\caption{2.5 DCF.}\end{figure}
\begin{lstlisting}[style=matlab]
dkr = 2*kmax/(Np-1); n_ang = 180;
area_proj = abs(kr_vec)*(2*pi/n_ang)*dkr;
[~,idx_dc]=min(abs(kr_vec));
area_proj(idx_dc) = pi*(dkr/2)^2 / n_ang;
area_proj2d = repmat(area_proj, 1, n_ang);
figure; plot(kr_vec, area_proj); xlabel('k_r'); ylabel('Area');
title('2.5 DCF (first angle)'); grid on; saveas(gcf, fullfile(plotDir,'2_5_dcf.png'));
\end{lstlisting}

\subsection{Direct Summation - Projection Data (2.6)}
Same row-by-row direct summation using the k-space data from the Projection-Slice Theorem. Images may be flipped/transposed; I corrected the orientation before comparison. \textbf{PSNR=20.30 dB, SSIM=0.4669, CPU=406.73 s.} Very slow due to large number of samples ($L=66060$).
\begin{figure}[H]\centering
\begin{subfigure}{0.48\textwidth}\includegraphics[width=\textwidth]{results/plots/2_6_image.png}\caption{DS.}\end{subfigure}\hfill
\begin{subfigure}{0.48\textwidth}\includegraphics[width=\textwidth]{results/plots/2_6_error.png}\caption{Error.}\end{subfigure}
\caption{2.6 Direct summation.}\end{figure}
\begin{figure}[H]\centering
\begin{subfigure}{0.48\textwidth}\includegraphics[width=\textwidth]{results/plots/2_6_horiz.png}\caption{Horizontal XS.}\end{subfigure}\hfill
\begin{subfigure}{0.48\textwidth}\includegraphics[width=\textwidth]{results/plots/2_6_vert.png}\caption{Vertical XS.}\end{subfigure}
\caption{2.6 Cross-sections.}\end{figure}
\begin{lstlisting}[style=matlab]
Np2=256; dxp=1; taup=Np2/2;
kxp=real(ktraj_proj(:)); kyp=imag(ktraj_proj(:));
dMp=area_proj2d(:).*kdata_proj(:);
np=(0:Np2-1)-taup;
ima_dsp=zeros(Np2,Np2);
t0=cputime;
for ix=1:Np2
    phx=exp(1j*2*pi*dxp*kxp*np(ix));
    phy=exp(1j*2*pi*dxp*kyp*np);
    ima_dsp(ix,:)=(dMp.*phx).'*phy;
    if mod(ix,64)==0, fprintf('  2.6 row %d/%d\n',ix,Np2); end
end
t_dsp=cputime-t0; clear phx phy;
ima_dsp=normImg(ima_dsp.');
% Auto-orient: try flips/transposes, pick best PSNR vs P
cands={ima_dsp,flipud(ima_dsp),fliplr(ima_dsp),ima_dsp.',...
    flipud(fliplr(ima_dsp)),rot90(ima_dsp,1),rot90(ima_dsp,3)};
bp=-Inf;
for ci=1:numel(cands)
    if all(size(cands{ci})==[256 256])
        pp=psnr(cands{ci},P); if pp>bp, bp=pp; ima_dsp=cands{ci}; end
    end
end
err_dsp=abs(ima_dsp-P);
psnr_dsp=psnr(ima_dsp,P); ssim_dsp=ssim(ima_dsp,P);
figure; imshow(ima_dsp,[0 1]); colormap gray; colorbar;
title('2.6 DS Radial'); saveas(gcf, fullfile(plotDir,'2_6_image.png'));
figure; imshow(err_dsp,[0 0.2]); colormap gray; colorbar;
title('2.6 Error'); saveas(gcf, fullfile(plotDir,'2_6_error.png'));
\end{lstlisting}

\subsection{2X Gridding - Projection Data (2.7)}
2X gridding with wg=4. To match the phantom orientation, I rotate the gridding output by $90^\circ$ counterclockwise before cropping and comparison. It achieved very similar quality to direct summation while being about 1000 times faster. \textbf{PSNR=20.30 dB, SSIM=0.4667, CPU=0.42 s.} This clearly shows the practical advantage of gridding.
\begin{figure}[H]\centering
\begin{subfigure}{0.32\textwidth}\includegraphics[width=\textwidth]{results/plots/2_7_full.png}\caption{Full output.}\end{subfigure}\hfill
\begin{subfigure}{0.32\textwidth}\includegraphics[width=\textwidth]{results/plots/2_7_crop.png}\caption{Cropped.}\end{subfigure}\hfill
\begin{subfigure}{0.32\textwidth}\includegraphics[width=\textwidth]{results/plots/2_7_error.png}\caption{Error.}\end{subfigure}
\caption{2.7 Gridding 2X.}\end{figure}
\begin{figure}[H]\centering
\begin{subfigure}{0.48\textwidth}\includegraphics[width=\textwidth]{results/plots/2_7_horiz.png}\caption{Horizontal XS.}\end{subfigure}\hfill
\begin{subfigure}{0.48\textwidth}\includegraphics[width=\textwidth]{results/plots/2_7_vert.png}\caption{Vertical XS.}\end{subfigure}
\caption{2.7 Cross-sections.}\end{figure}
\begin{lstlisting}[style=matlab]
t0=cputime;
ima_g2p_full=gridkb(kdata_proj,ktraj_proj,area_proj2d,256,2,4,'image');
t_g2p=cputime-t0;
ima_g2p_full=normImg(rot90(ima_g2p_full,1));
c5=256; h5=128;
ima_g2p=normImg(ima_g2p_full(c5-h5+1:c5+h5,c5-h5+1:c5+h5));
err_g2p=abs(ima_g2p-P);
psnr_g2p=psnr(ima_g2p,P); ssim_g2p=ssim(ima_g2p,P);
figure; imshow(ima_g2p_full,[0 1]); colormap gray; colorbar;
title('2.7 Full Output'); saveas(gcf, fullfile(plotDir,'2_7_full.png'));
figure; imshow(ima_g2p,[0 1]); colormap gray; colorbar;
title('2.7 Cropped 256x256'); saveas(gcf, fullfile(plotDir,'2_7_crop.png'));
figure; imshow(err_g2p,[0 0.2]); colormap gray; colorbar;
title('2.7 Error'); saveas(gcf, fullfile(plotDir,'2_7_error.png'));
figure; plot(ima_g2p(128,:),'b'); hold on; plot(P(128,:),'r--');
legend('Grid2X','Ref'); title('2.7 Horiz XS'); grid on; saveas(gcf,fullfile(plotDir,'2_7_horiz.png'));
figure; plot(ima_g2p(:,128),'b'); hold on; plot(P(:,128),'r--');
legend('Grid2X','Ref'); title('2.7 Vert XS'); grid on; saveas(gcf,fullfile(plotDir,'2_7_vert.png'));
\end{lstlisting}

\section{Part III: Gridding for Radial Data}

\subsection{K-Space Trajectory (3.1)}
Loaded \texttt{shepplogan\_radial\_data.mat} ($129\times300$). Unlike the projection data, these radial lines start at the center and extend outward in one direction only (center-out, one-sided).
\begin{figure}[H]\centering\includegraphics[width=0.5\textwidth]{results/plots/3_1_trajectory.png}\caption{3.1 One-sided radial trajectory.}\end{figure}
\begin{lstlisting}[style=matlab]
load('shepplogan_radial_data.mat');
figure; plot(real(ktraj(:)),imag(ktraj(:)),'b.','MarkerSize',1);
xlabel('k_x'); ylabel('k_y'); title('3.1 Radial Trajectory (one-sided)');
axis equal; grid on; saveas(gcf, fullfile(plotDir,'3_1_trajectory.png'));
\end{lstlisting}

\subsection{Geometric DCF (3.2)}
Same geometric formula as 2.5 but with $k_r \in [0, k_{max}]$ instead of $[-k_{max}, k_{max}]$.
\begin{figure}[H]\centering\includegraphics[width=0.5\textwidth]{results/plots/3_2_dcf.png}\caption{3.2 One-sided radial DCF.}\end{figure}
\begin{lstlisting}[style=matlab]
[Ns,Nl]=size(ktraj); kmax_r=max(abs(ktraj(:)));
kr_r=abs(ktraj(:,1)); dkr_r=kmax_r/(Ns-1);
area_rad1d=kr_r*(2*pi/Nl)*dkr_r;
area_rad1d(1)=pi*(dkr_r/2)^2/Nl;
area_radial=repmat(area_rad1d,1,Nl);
figure; plot(kr_r,area_rad1d); xlabel('k_r'); ylabel('Area');
title('3.2 DCF (one radial line)'); grid on; saveas(gcf, fullfile(plotDir,'3_2_dcf.png'));
\end{lstlisting}

\subsection{Direct Summation (3.3)}
Row-by-row direct summation on the one-sided radial data. Auto-oriented. \textbf{PSNR=18.23 dB, SSIM=0.4635, CPU=274.14 s.} Visible radial streak artifacts. Similar quality to the projection data case.
\begin{figure}[H]\centering
\begin{subfigure}{0.48\textwidth}\includegraphics[width=\textwidth]{results/plots/3_3_image.png}\caption{DS.}\end{subfigure}\hfill
\begin{subfigure}{0.48\textwidth}\includegraphics[width=\textwidth]{results/plots/3_3_error.png}\caption{Error.}\end{subfigure}
\caption{3.3 Direct summation.}\end{figure}
\begin{figure}[H]\centering
\begin{subfigure}{0.48\textwidth}\includegraphics[width=\textwidth]{results/plots/3_3_horiz.png}\caption{Horizontal XS.}\end{subfigure}\hfill
\begin{subfigure}{0.48\textwidth}\includegraphics[width=\textwidth]{results/plots/3_3_vert.png}\caption{Vertical XS.}\end{subfigure}
\caption{3.3 Cross-sections.}\end{figure}
\begin{lstlisting}[style=matlab]
Nr=256; dxr=1; taur=Nr/2;
kxr=real(ktraj(:)); kyr=imag(ktraj(:));
dMr=area_radial(:).*kdata(:);
nr=(0:Nr-1)-taur;
ima_dsr=zeros(Nr,Nr);
t0=cputime;
for ix=1:Nr
    phx=exp(1j*2*pi*dxr*kxr*nr(ix));
    phy=exp(1j*2*pi*dxr*kyr*nr);
    ima_dsr(ix,:)=(dMr.*phx).'*phy;
    if mod(ix,64)==0, fprintf('  3.3 row %d/%d\n',ix,Nr); end
end
t_dsr=cputime-t0; clear phx phy;
ima_dsr=normImg(ima_dsr.');
% Auto-orient
cands={ima_dsr,flipud(ima_dsr),fliplr(ima_dsr),ima_dsr.',...
    flipud(fliplr(ima_dsr)),rot90(ima_dsr,1),rot90(ima_dsr,3)};
bp=-Inf;
for ci=1:numel(cands)
    if all(size(cands{ci})==[256 256])
        pp=psnr(cands{ci},P); if pp>bp, bp=pp; ima_dsr=cands{ci}; end
    end
end
err_dsr=abs(ima_dsr-P);
psnr_dsr=psnr(ima_dsr,P); ssim_dsr=ssim(ima_dsr,P);
figure; imshow(ima_dsr,[0 1]); colormap gray; colorbar;
title('3.3 DS Radial'); saveas(gcf, fullfile(plotDir,'3_3_image.png'));
figure; imshow(err_dsr,[0 0.2]); colormap gray; colorbar;
title('3.3 Error'); saveas(gcf, fullfile(plotDir,'3_3_error.png'));
figure; plot(ima_dsr(128,:),'b'); hold on; plot(P(128,:),'r--');
legend('DS','Ref'); title('3.3 Horiz XS'); grid on; saveas(gcf,fullfile(plotDir,'3_3_horiz.png'));
figure; plot(ima_dsr(:,128),'b'); hold on; plot(P(:,128),'r--');
legend('DS','Ref'); title('3.3 Vert XS'); grid on; saveas(gcf,fullfile(plotDir,'3_3_vert.png'));
\end{lstlisting}

\subsection{2X Gridding (3.4)}
To match the phantom orientation, I rotate the gridding output by $90^\circ$ counterclockwise before cropping and comparison. The result has very similar quality to direct summation but is about 900 times faster. \textbf{PSNR=18.23 dB, SSIM=0.4636, CPU=0.30 s.}
\begin{figure}[H]\centering
\begin{subfigure}{0.32\textwidth}\includegraphics[width=\textwidth]{results/plots/3_4_full.png}\caption{Full output.}\end{subfigure}\hfill
\begin{subfigure}{0.32\textwidth}\includegraphics[width=\textwidth]{results/plots/3_4_crop.png}\caption{Cropped.}\end{subfigure}\hfill
\begin{subfigure}{0.32\textwidth}\includegraphics[width=\textwidth]{results/plots/3_4_error.png}\caption{Error.}\end{subfigure}
\caption{3.4 Gridding 2X.}\end{figure}
\begin{figure}[H]\centering
\begin{subfigure}{0.48\textwidth}\includegraphics[width=\textwidth]{results/plots/3_4_horiz.png}\caption{Horizontal XS.}\end{subfigure}\hfill
\begin{subfigure}{0.48\textwidth}\includegraphics[width=\textwidth]{results/plots/3_4_vert.png}\caption{Vertical XS.}\end{subfigure}
\caption{3.4 Cross-sections.}\end{figure}
\begin{lstlisting}[style=matlab]
t0=cputime;
ima_g2r_full=gridkb(kdata,ktraj,area_radial,256,2,4,'image');
t_g2r=cputime-t0;
ima_g2r_full=normImg(rot90(ima_g2r_full,1));
c6=256;
ima_g2r=normImg(ima_g2r_full(c6-h5+1:c6+h5,c6-h5+1:c6+h5));
err_g2r=abs(ima_g2r-P);
psnr_g2r=psnr(ima_g2r,P); ssim_g2r=ssim(ima_g2r,P);
figure; imshow(ima_g2r_full,[0 1]); colormap gray; colorbar;
title('3.4 Full Output'); saveas(gcf, fullfile(plotDir,'3_4_full.png'));
figure; imshow(ima_g2r,[0 1]); colormap gray; colorbar;
title('3.4 Cropped 256x256'); saveas(gcf, fullfile(plotDir,'3_4_crop.png'));
figure; imshow(err_g2r,[0 0.2]); colormap gray; colorbar;
title('3.4 Error'); saveas(gcf, fullfile(plotDir,'3_4_error.png'));
figure; plot(ima_g2r(128,:),'b'); hold on; plot(P(128,:),'r--');
legend('Grid2X','Ref'); title('3.4 Horiz XS'); grid on; saveas(gcf,fullfile(plotDir,'3_4_horiz.png'));
figure; plot(ima_g2r(:,128),'b'); hold on; plot(P(:,128),'r--');
legend('Grid2X','Ref'); title('3.4 Vert XS'); grid on; saveas(gcf,fullfile(plotDir,'3_4_vert.png'));
\end{lstlisting}

\section{Results Summary}
\begin{table}[H]\centering\caption{All results.}
\begin{tabular}{l r r r}\toprule
\textbf{Method} & \textbf{PSNR (dB)} & \textbf{SSIM} & \textbf{CPU (s)} \\\midrule
1.4 Direct Sum (ref)     & $\infty$ & 1.0000 & 0.98 \\
1.5 No DCF               & 10.86 & 0.2729 & 0.95 \\
1.6 Double FOV           & $\infty$ & 1.0000 & 2.36 \\
1.7 Half Pixel           & 33.67 & 0.9150 & 3.62 \\
1.8 Double Pixel         & 18.95 & 0.6160 & 0.36 \\
1.9 Gridding 1X          & 23.91 & 0.7537 & 0.05 \\
1.10 Gridding 2X         & 44.35 & 0.9906 & 0.09 \\
1.11 ScatInterp 1X       & 22.85 & 0.6705 & 0.06 \\
1.12 ScatInterp 2X       & 24.33 & 0.7350 & 0.27 \\\midrule
2.2 FBP                  & 26.82 & 0.6799 & 0.08 \\
2.3 Naive BP             &  4.72 & 0.1535 & 0.08 \\
2.6 DS (Projection)      & 20.30 & 0.4669 & 406.73 \\
2.7 Grid (Projection)    & 20.30 & 0.4667 & 0.42 \\\midrule
3.3 DS (Radial)          & 18.23 & 0.4635 & 274.14 \\
3.4 Grid (Radial)        & 18.23 & 0.4636 & 0.30 \\\bottomrule
\end{tabular}\end{table}

\section{Conclusion}
In this homework, I learned that non-Cartesian MRI reconstruction needs careful handling of sampling density, and that the choice of reconstruction method involves a trade-off between speed and accuracy.
\begin{enumerate}
  \item \textbf{Density compensation} is essential for correcting non-uniform sampling density. Without it, the image quality degrades significantly (PSNR dropped from $\infty$ to 10.86 dB).
  \item \textbf{Ramp filtering} is necessary for backprojection to avoid $1/r$ blurring. Naive backprojection without the filter gave only 4.72 dB PSNR.
  \item \textbf{Gridding reconstruction} provides an excellent trade-off between accuracy and speed. In the projection and radial cases here, it achieved nearly the same PSNR/SSIM as direct summation at roughly three orders of magnitude lower CPU time.
\end{enumerate}
The choice of pixel size and FOV affects both quality and computation time. Overall, gridding with 2X oversampling is a strong practical choice for non-Cartesian MRI reconstruction in this homework.


\end{document}
